\chapter{Autoridad certificadora y sellado de tiempo}\label{cap:ca}

Qadra incluye una autoridad certificadora (CA) interna y una autoridad de
sellado de tiempo (TSA) local, ambas construidas sobre OpenSSL. Este capítulo
resume su estructura y su operación; el detalle de cada paso está en
\archivo{docs/ca/operations.md}, \archivo{pki/README.md} y el Anexo D del
reporte técnico.

\importante{La CA y la TSA son internas y de prototipo. No son un prestador
de servicios de certificación, no emiten e.firma ni FIREL y no acreditan la
identidad civil del titular. Un certificado emitido prueba únicamente que su
titular posee la clave privada correspondiente.}

% ---------------------------------------------------------------------
\section{Estructura de la PKI}\label{sec:ca-estructura}

La PKI tiene un solo nivel: una raíz autofirmada que emite directamente los
certificados de entidad final. No hay autoridades intermedias ni OCSP; la
revocación se publica mediante una lista de revocación (CRL). Todas las
claves son RSA de 3072 bits y todas las firmas usan SHA-256. Los perfiles se
definen en \archivo{pki/openssl.cnf} y \archivo{pki/tsa.cnf}
(\cref{tab:ca-perfiles}).

\begin{table}[H]
  \centering\small
  \caption{Perfiles de certificado de la PKI interna.}
  \label{tab:ca-perfiles}
  \begin{tabularx}{\textwidth}{L{3.4cm}L{4.9cm}Y}
    \toprule
    \textbf{Perfil} & \textbf{Sección OpenSSL} & \textbf{Usos y vigencia} \\
    \midrule
    Raíz & \texttt{v3\_root\_ca} &
      \texttt{CA:true}, \texttt{keyCertSign} y \texttt{cRLSign}; 1825 días. \\
    Firma documental (\texttt{document-signing}) &
      \texttt{v3\_internal\_declaration} &
      \texttt{digitalSignature} y \texttt{nonRepudiation}, sin usos
      extendidos; 365 días. \\
    Acceso del Administrador (\texttt{partner}) & \texttt{v3\_end\_entity} &
      Los mismos usos más \texttt{clientAuth} y \texttt{emailProtection};
      365 días. \\
    Sellado de tiempo (\texttt{time-stamping}) & \texttt{v3\_tsa} &
      \texttt{digitalSignature} y uso extendido crítico
      \texttt{timeStamping}; 365 días. \\
    \bottomrule
  \end{tabularx}
\end{table}

La política \texttt{policy\_internal} exige que el país (C) y la organización
(O) del sujeto coincidan con los de la raíz y que exista un nombre común (CN).
Las extensiones solicitadas en el CSR no se copian. La CRL vence a los 7 días
de generada, por lo que debe renovarse periódicamente
(\cref{sec:ca-revocacion}).

El material vive fuera del repositorio, en dos directorios de trabajo:
\texttt{PKI\_CA\_DIR} para la CA (por omisión \archivo{pki-ca} bajo el
directorio actual) y \texttt{TSA\_DIR} para la TSA (por omisión
\archivo{pki-tsa} junto al anterior). En una instalación desplegada
corresponden a \archivo{data/ca} y \archivo{data/tsa} bajo la raíz de la
instalación. Ambos contienen claves privadas y nunca se versionan.

\section{Creación de la CA y de la TSA}\label{sec:ca-creacion}

Los guiones de \archivo{pki/} se invocan directamente o mediante la CLI
(\cref{cap:cli}). La creación se hace una sola vez por instalación:

\begin{verbatim}
export PKI_CA_DIR="$HOME/pki-ca"
bash pki/init-ca.sh          # raíz, índice y contadores
bash pki/issue-tsa-cert.sh   # clave y certificado de la TSA
bash pki/gen-crl.sh          # primera CRL
\end{verbatim}

\archivo{init-ca.sh} no sobrescribe una CA existente. La TSA responde
solicitudes RFC 3161 con \texttt{openssl ts} según \archivo{pki/tsa.cnf}; el
subcomando \texttt{timestamp --mock} de la CLI la utiliza y falla si el
directorio de la TSA no existe. El certificado de la TSA aparece en el inventario
del portal solo para consulta y descarga: no se emite ni se revoca desde ahí.

\nota{Instalar el portal no crea una autoridad nueva. El servicio adopta la
raíz, el índice, los contadores y la TSA existentes, y fija la huella de la
raíz en su configuración.}

\section{Servicio y portal de la CA}\label{sec:ca-portal}

La CA se opera desde un servicio propio, separado del despacho:

\begin{itemize}
  \item \textbf{Binario \texttt{qadra-ca}} (\archivo{crates/bin/src/ca/}).
        Expone el acceso del operador bajo \texttt{/api/v1/auth}, las
        operaciones en \texttt{POST /api/v1/ca/operations} y
        \texttt{/healthz}. Escucha solo en \emph{loopback}. El contrato HTTP
        está en \archivo{docs/ca/api.md}.
  \item \textbf{Portal web} (\archivo{web/apps/ca}). Se construye aparte de la
        aplicación del despacho con \texttt{npm run build:ca} desde
        \archivo{web/}. En desarrollo usa el puerto 18088 y dirige
        \texttt{/api} a \texttt{127.0.0.1:18089}; el destino se cambia con
        \texttt{CA\_API\_PROXY\_TARGET}.
  \item \textbf{Ejecutor de operaciones} (\archivo{ops/ca/command.py}), que
        invoca los guiones PKI y el mantenimiento de CRL existentes.
\end{itemize}

El servicio comparte PostgreSQL y Redis con el despacho, pero con identidad
propia: el esquema \texttt{qadra\_ca} (tablas \texttt{users} y
\texttt{audit\_events}) y un rol de ejecución con permisos mínimos y
\texttt{search\_path=qadra\_ca,pg\_temp}. Las credenciales, las sesiones y la
auditoría de la CA no se mezclan con las del despacho; una cuenta de
Administrador del despacho no da acceso al portal. Las unidades
\texttt{qadra-ca} y \texttt{qadra-ca-web} las genera
\archivo{ops/ca/host.py} (\cref{cap:despliegue}).

\section{Alta del operador}\label{sec:ca-alta}

La CA admite un único operador, fijado por UUID y correo en la configuración.
Su alta se hace por terminal, una sola vez, después de preparar el esquema
(\archivo{crates/infrastructure/src/ca_identity/schema.sql}) y el rol:

\begin{enumerate}
  \item Prepare un archivo privado de alta (modo \texttt{0600}, sin enlaces
        simbólicos) con la base, el correo del operador, la KEK de la CA de
        32 bytes en Base64 y las conexiones administrativa y de ejecución.
        Los campos exactos están en \archivo{docs/ca/operations.md}.
  \item Ejecute el alta:
\begin{verbatim}
qadra-ca enroll-operator \
  --setup-config /ruta/privada/ca-setup.json \
  --enrollment-out /ruta/privada/ca-enrollment.json
\end{verbatim}
  \item Introduzca la contraseña (12 a 1024 bytes) dos veces. El comando
        escribe en el archivo de salida el secreto TOTP, su URI y ocho
        códigos de recuperación.
  \item Registre el secreto en una aplicación de autenticación e introduzca
        el código de seis dígitos. Solo entonces se crean la cuenta y el
        evento de auditoría \texttt{ca.operator\_enrolled}.
  \item Cree la configuración del servicio con el UUID del recibo, el mismo
        correo y la misma KEK, y arranque:
        \texttt{qadra-ca --config /ruta/privada/ca-service.json}.
\end{enumerate}

\importante{El alta no se repite para reemplazar una cuenta existente. La
KEK protege el secreto TOTP almacenado: si se pierde o se cambia, el segundo
factor del operador deja de ser recuperable. Conserve el archivo de alta, la
KEK y las URL de conexión fuera del repositorio, de los registros y de las
capturas.}

La sección opcional \texttt{operations} de la configuración activa las
operaciones PKI; sin ella el servicio solo ofrece el acceso. En ella se
declaran las rutas absolutas del intérprete, el ejecutor, la raíz de la
instalación, la CLI y los guiones, además de \texttt{expected\_root}, la
huella SHA-256 de la raíz adoptada.

\section{Solicitud, emisión y vínculo}\label{sec:ca-emision}

El titular genera su clave y su solicitud (CSR) en su propio equipo; la CA
recibe únicamente el CSR público. La clave puede generarse desde la aplicación
(Manual de usuario, capítulo \enquote{Mi firma digital}) o con OpenSSL,
con los comandos de \archivo{pki/README.md}. En el portal, el operador
sigue este recorrido:

\begin{enumerate}
  \item En \textbf{Solicitudes}, pega el CSR en formato PEM (máximo 16 KiB) y
        elige el perfil \textbf{Firma documental} o \textbf{Partner}.
  \item Completa \textbf{Titular revisado}, \textbf{Referencia de la
        solicitud} y \textbf{Criterio de emisión}, y pulsa \textbf{Revisar
        solicitud}. Contrasta el sujeto y las huellas del CSR, de la clave
        pública y de la CA con los datos que entregó el titular.
  \item Pulsa \textbf{Confirmar y emitir} y conserva el recibo con la serie y
        la huella. El certificado se descarga desde \textbf{Certificados}.
\end{enumerate}

El CSR debe usar RSA de 3072 bits, exponente 65537 y firma SHA-256. Una
solicitud cuyo sujeto no cumple la política se rechaza antes de reservar una
serie. Si la respuesta se pierde, use \textbf{Retomar esta emisión} o
\textbf{Historial}; no envíe la solicitud como una operación nueva. Sin
portal, la misma emisión se obtiene con
\texttt{despacho-cli pki issue-request --csr \ldots{} --name \ldots{}} y,
para el acceso del Administrador, \texttt{--purpose partner}.

\nota{Emitir un certificado no crea una cuenta ni lo vincula. El titular
registra el certificado en Qadra y firma una prueba de posesión; ese vínculo
es el que habilita la firma individual o el acceso por certificado. Retirar
el vínculo y revocar el certificado son operaciones distintas. Véanse
\archivo{docs/individual-signing-guide.md} y
\archivo{docs/owner-certificate-login-operations.md}.}

\section{Revocación y publicación de la CRL}\label{sec:ca-revocacion}

La revocación tiene dos pasos, y el primero no surte efecto por sí solo:

\begin{enumerate}
  \item \textbf{Revocación local.} En \textbf{Certificados}, elija
        \textbf{Revocar certificado}, confirme la serie y la huella, y
        escriba el motivo. Esto solo marca el índice de la CA; el inventario
        muestra \enquote{Pendiente de publicar en Qadra}.
  \item \textbf{Publicación.} En \textbf{Publicación}, revise la revisión
        vigente, confirme la ventana de mantenimiento y pulse \textbf{Renovar
        y publicar CRL}. El procedimiento prepara la CRL candidata, detiene
        la API y la web del despacho, toma un respaldo, publica la lista
        como una nueva revisión de confianza en PostgreSQL y reabre el
        servicio.
\end{enumerate}

Qadra decide con la confianza publicada en la base de datos (raíz, CRL y
número de revisión), no con los archivos locales de la CA. Mientras la
publicación no termine, el certificado revocado sigue siendo aceptado. Una
operación pendiente bloquea otras mutaciones: retómela desde
\textbf{Historial} en lugar de iniciar otra.

La renovación también puede ejecutarse sin portal, a mano o mediante el
temporizador diario, con \archivo{ops/deploy/renew_crl.py}
(\cref{sec:op-crl}). En un entorno local, la confianza se publica con
\texttt{despacho-cli credential-trust publish} y una conexión administrativa
(\archivo{pki/README.md}, \enquote{Publicar confianza para declaraciones}).
La primera publicación fija la raíz de la instalación; las siguientes solo
admiten una CRL con número mayor y no sustituyen la raíz.

\section{Raíces fijadas en el validador}\label{sec:ca-raices}

El validador de evidencia del navegador (\cref{cap:verificacion}) no confía
en una raíz por el hecho de venir dentro del paquete. Solo admite las raíces
cuya huella SHA-256 consta en
\archivo{web/apps/qadra/src/lib/evidence/pinned-roots.json}, que se integra
en la aplicación al construirla. Cada entrada tiene exactamente tres campos:
\texttt{fingerprint} (64 dígitos hexadecimales en minúsculas),
\texttt{purposes} (uno o más de \texttt{server-seal},
\texttt{individual-signing} y \texttt{time-stamping}) y \texttt{label}.

Para construir la web contra otra CA, defina \texttt{QADRA\_PINNED\_ROOTS}
con la ruta absoluta de un archivo del mismo formato;
\archivo{web/apps/qadra/astro.config.mjs} lo valida y detiene la
construcción si es inválido. \archivo{scripts/web-demo.sh} fija así su CA
desechable. Además, el validador consulta
\texttt{GET /api/v1/evidence/trust} para conocer la confianza que el servidor
admite en ese momento y reporta cualquier diferencia con el conjunto fijado.

\importante{Si se sustituye la raíz de una instalación, agregue la huella
nueva y conserve la anterior en el archivo. Al retirarla, los paquetes de
evidencia emitidos bajo esa raíz dejarían de validarse.}

\section{Respaldo y restauración}\label{sec:ca-restauracion}

No existe un comando único de restauración de la CA; el procedimiento
administrativo está en \archivo{docs/ca/restore.md}. Sus puntos esenciales
son:

\begin{itemize}
  \item \textbf{Qué conservar del mismo punto de recuperación:} el volcado de
        la base (incluido el esquema \texttt{qadra\_ca}), la configuración
        privada del servicio con su KEK, todo \archivo{data/ca} (raíz, clave,
        índice, contadores, certificados y los directorios
        \archivo{operations}, \archivo{revocations} y
        \archivo{publications}), \archivo{data/tsa} y el estado de Redis.
  \item \textbf{Orden:} detener y enmascarar el servicio de la CA y los del
        despacho; restaurar base, roles y archivos; invalidar las sesiones en
        Redis; comprobar el acceso del operador y la confianza publicada
        antes de reabrir.
  \item \textbf{Qué no hacer:} repetir \texttt{enroll-operator}, generar otra
        KEK, regenerar la raíz o la TSA, bajar contadores o borrar registros
        de operaciones para destrabar un arranque.
\end{itemize}

Una publicación de CRL que quedó pendiente al tomar el respaldo se retoma
con su identificador original desde \textbf{Historial}; no se genera otra
lista. El respaldo y la restauración generales de la instalación se
describen en las \cref{sec:op-respaldo,sec:op-restauracion}.
